Agronomic management of the crop and its nitrogen supply is the principal
in-field lever on drainage nitrate, complementary to the drainage-system and
edge-of-field practices treated elsewhere. A large plot-, field-,
watershed-, and regional-scale literature converges on a consistent picture:
drainage nitrate concentration and load are governed jointly by nitrogen (N)
input, crop and rotation choice, the amount and seasonal timing of drainage
flow, and soil processing of N, so that no single practice is universally
decisive but several act in predictable directions
\parencite{Dinnes2002}. Under conventional maize and maize--soybean systems
flow-weighted NO$_3^-$-N routinely exceeds the $10$~mg~L$^{-1}$
drinking-water standard---more than $70\%$ of site-years in one twelve-state
synthesis \parencite{Liu2021}, and corn, corn--soybean, and soybean
site-years averaging $14.0$, $13.5$, and $12.1$~mg~L$^{-1}$ across a
continental database \parencite{Hertzberger2019}. The underlying rate--loss
and hydrologic mechanisms are developed in the nitrate-export section; here
the focus is on how rate, source, timing, cover crops, rotation, and tillage
alter that loss and at what cost.

\subsection{Nitrogen Rate, Source, and Residual Soil Nitrate}

Drainage nitrate loss rising with N application rate is the most consistently
documented in-field relationship, and it is paired with a yield plateau near
the economic optimum. In Iowa production fields, yearly nitrate mass loss was
greater for high- than for medium- and low-N treatments, with concentrations
of $5$--$45$~mg~N~L$^{-1}$ \parencite{Jaynes2001}; four-year losses ranged
$22$--$68$~kg~N~ha$^{-1}$ and rose with rate in both continuous corn and
corn--soybean \parencite{Helmers2012b}. Increasing N above the agronomic
optimum yields diminishing returns: at $168$--$252$~kg~N~ha$^{-1}$ in
Illinois maize yield differed little while cumulative loads were
statistically similar \parencite{Li2023b}, echoing early Iowa work that
identified a single $150$~kg~ha$^{-1}$ application as a favorable compromise
\parencite{Azevedo1997}; a $10\%$ rate cut lowered yield without significantly
reducing loss at one site, underscoring site and hydrologic dependence
\parencite{Gentry2024}. Regional syntheses formalize the pattern: maize yield
followed a quadratic response peaking near $193$~kg~N~ha$^{-1}$ while drainage
N loss increased exponentially with rate, and yield-scaled loss rose sharply
once N surplus exceeded $\sim23$~kg~N~ha$^{-1}$, identifying zero surplus as a
practical target \parencite{Zhao2016}. N rate was the strongest predictor of
annual concentration ($R^2=0.59$) with no yield gain above
$150$--$224$~kg~ha$^{-1}$ \parencite{Hertzberger2019}, the most cost-effective
high-rate scenarios also produced the highest losses \parencite{Liu2021}, and
across syntheses each unit rise in rate produced $\sim0.12$~kg~NO$_3^-$-N
loss, rate explaining $27\%$ of variability
\parencite{Chatterjee2020b,Christianson2015b,Ritter2015}.

Residual soil nitrate (RSN) links rate to leaching risk and creates
legacies. RSN was minimally affected below the economic optimum N rate (EONR)
and increased sharply only above it \parencite{Bandura2025}; continuous corn
accumulated up to $225$~kg~ha$^{-1}$ RSN in dry years, flushed to drainage in
subsequent wet years \parencite{Randall2001}, and a very high fall rate
continued to affect nitrate and residual N for more than four years
\parencite{Jaynes2015}. Encouragingly, corn yields rose over three decades at
nearly constant rate while drainage nitrate stayed flat, indicating improved
use efficiency but no loss reduction from yield gains alone
\parencite{Clifford2023}; the leaching breakpoint lies above the agronomic
optimum, with a larger buffer in rotations \parencite{Pasley2021}.

Nitrogen source matters, with manure generally increasing loss relative to
inorganic fertilizer at comparable rates. Liquid swine manure produced
greater concentrations and losses than UAN with no yield difference
\parencite{Bakhsh2005}, liquid-hog manure raised soil and drainage nitrate
above inorganic fertilizer with equivalent yields \parencite{Elmi2005}, and
annually repeated fall manure produced the highest losses, timing of manure
addition controlling loss more than rotation \parencite{HernandezRamirez2011}.
Applying manure before both crops rather than only corn raised loads
\parencite{Dougherty2020}, and fall manure leached more than spring manure
\parencite{Bakhsh2007,Rekha2011}. Organic sources generally boosted yield but
increased losses \parencite{Liu2021}, though one meta-analysis found no
significantly greater dissolved N load \parencite{Christianson2015b}.

\subsection{Nitrogen Timing, Placement, and Enhanced-Efficiency Products}

Shifting N from fall to spring or in-season application generally lowers
drainage nitrate, but effects are smaller and less consistent than rate or
cropping-system changes. In Minnesota rotations fall N without nitrapyrin
produced the highest losses, with reductions relative to fall of
$13$--$18\%$ (spring, split, fall+nitrapyrin) \parencite{Randall2003,Randall2005};
fall N produced greater concentration and load than spring N in Illinois
\parencite{Gentry2024}, and spring or side-dress N cut concentrations while
maintaining yield \parencite{Pittelkow2017,Vetsch2019}. Minnesota work found
spring application raised use efficiency $>20\%$ and reduced tile losses
$\sim36\%$ relative to fall \parencite{Dinnes2002}, though the fall-versus-spring
effect was often smaller than weather, drainage volume, and rate
\parencite{Liu2021}, and out-of-season and pre-plant timings had higher
concentrations than at-plant and side-dress timings \parencite{Hertzberger2019}.

Split, side-dress, and placement effects are mixed. Split fertilization
reduced export only at lower rates \parencite{Liu2021} and was not
significantly different from non-split in a continental synthesis
\parencite{Hertzberger2019}; side-dress at V2--V12 gave no consistent
difference \parencite{Jaynes2013}, and midseason splits sometimes raised
concentration and RSN \parencite{Jaynes2006,Bandura2025}. Split-applied urea
modestly improved water quality \parencite{Waring2022}, and remote-sensing
variable-rate side-dress averaged a $40\%$ reduction over 15 years, mainly via
$33\%$ lower input and avoidance of peak flow, while maintaining yield
\parencite{Wilson2020b}. Placement was inconsistent: injection and
incorporation sometimes had higher concentrations than surface application
\parencite{Hertzberger2019}, and a localized-compaction applicator gave $16\%$
lower concentrations than a spoke injector \parencite{Bakhsh2010}. Among
enhanced-efficiency products, nitrapyrin generally did not reduce nitrate or
improve yield alone \parencite{Vetsch2019,Pittelkow2017} but reduced fall-N
losses $10$--$18\%$ in Minnesota rotations \parencite{Randall2003,Randall2005};
polymer-coated products were more robust, with ESN lowering combined N$_2$O,
NH$_3$, and NO$_3^-$ losses $47$--$49\%$ below urea while maintaining yield,
the largest benefit in wet periods \parencite{Menegaz2025,Nash2015,Nash2015b}.
Reduced-rate summer fertigation cut concentration and load $\sim16\%$ without
reducing silage yield \parencite{Gamble2018}, and fall anhydrous ammonia with
nitrapyrin matched spring application environmentally \parencite{Kaur2024,Kaur2024c}.

\subsection{Cover Crops: Reductions, Mechanisms, and Yield}

Cover crops---particularly fall-planted winter cereal rye---are among the
most consistent and largest in-field reducers of drainage nitrate. The
clearest, most repeatable signal is reduced concentration. In a long-term
Iowa study rye cut concentrations $48\%$ over five years and oat $26\%$,
though neither significantly reduced cumulative loads because annual drainage
varied \parencite{Kaspar2012}; rye reduced concentrations $>30\%$ even at low
biomass \parencite{Waring2020}, and reviews report reductions of
$13$--$94\%$ \parencite{Ward2018}. Where rye is well established, large load
reductions follow: split N plus rye reduced flow-weighted concentration and
load $35\%$ and $37\%$ without yield penalty \parencite{PrezaFontes2021}; rye
added to spring and fall systems cut load $40\%$ and $47\%$, reaching
$61$--$64\%$ in a wet year \parencite{Ruffatti2019}; a no-till rye system
lowered losses $59$--$67\%$ \parencite{OBrien2022}; rye contributed a $43\%$
decrease in cumulative four-year load \parencite{Gentry2024}; and Indiana
watershed monitoring found field-scale reductions of $58$--$99\%$, pooling to
$65\%$ \parencite{Vincent2025}. Multi-year and northern-margin studies show
meaningful but often statistically marginal cumulative effects: rye reduced
cumulative drainage $12.1\%$ and N loss $20.4\%$ over six Iowa years, with
significant reductions only in favorable years \parencite{MartinezFeria2016};
$11\%$ discharge and $13\%$ loss reductions in Minnesota \parencite{Strock2004};
and in a manured silage system rye lowered concentration from $53$ to
$39$~mg~L$^{-1}$ but loads remained high, warranting additional practices
\parencite{Krueger2013}.

The mechanism is assimilation of soil mineral N into biomass during the
fall--spring fallow, shifting nitrate out of leaching-vulnerable depths.
Cereal rye and radish accumulated $130$--$250$~kg~N~ha$^{-1}$ and moved
NO$_3$-N from environmental (20--80~cm) to agronomic depths, high-C/N
overwintering rye best protecting N in wet springs \parencite{Lacey2015};
living cover consistently lowered soil and vadose-water nitrate
\parencite{Bosompemaa2021,Strock2004,Qi2011b}. Effectiveness is therefore
strongly biomass- and weather-dependent: biomass and reductions were weak
under poor conditions \parencite{PrezaFontes2021,OBrien2022}, cereal rye
biomass varied $1.4$--$5.9$~Mg~ha$^{-1}$ (highest under late termination)
\parencite{Gentry2025,Birru2023}, and success depended on early establishment,
achieved in only one of four years after corn in one Minnesota study
\parencite{Strock2004,Bawa2023}. Regional modeling explained $>95\%$ of
site-to-site variation, rye being more effective in warmer, less-wet southern
locations \parencite{Malone2014b}. Effects on drainage \emph{volume} are
modest and often non-significant, so load reductions are driven mainly by
concentration: rye cut May--July drainage $21\%$ via enhanced spring
evapotranspiration but annual drainage only $9\%$ (non-significant)
\parencite{Qi2010b,Qi2011b}, and in the 2012 drought rye did not deplete soil
water at the cash crop's expense \parencite{Daigh2014b}.

Rye cover cropping generally imposes no consistent yield penalty, especially
for soybean \parencite{Strock2004,Waring2020,Qi2011b}, with projected-climate
differences of only $\sim0.1$~Mg~ha$^{-1}$ and no continuous-corn penalty in
simulation \parencite{Malone2020,Birru2023}. Documented exceptions concentrate
in maize, drought years, and silage---a $6\%$ average maize reduction (34\% in
2012) \parencite{MartinezFeria2016} and $16\%$ lower corn silage after manured
rye \parencite{Krueger2013}. Management modulates outcomes: earlier fall
sowing is consistently more effective (mid-September sowing reduced losses
$45\%$ versus October dates) \parencite{Feyereisen2006b}, and delayed spring
termination raises biomass with limited yield cost \parencite{Birru2023,Chatterjee2025}.
In harvested-rye systems, fertilizing the rye boosts biomass and energy with
only a small N-loss cost \parencite{Malone2018,Malone2023}. Synthesizing
across studies, well-established rye clusters around $40$--$60\%$ N-loss
reduction with $25$--$50\%$ concentration reduction
\parencite{Malone2018,Chatterjee2025,Malone2014b,Malone2023}; Iowa
climate-change projections had rye reduce annual loads from $47.5$ to
$18.9$~kg~N~ha$^{-1}$ ($60\%$) under baseline and $63.6$ to
$17.0$~kg~N~ha$^{-1}$ ($73\%$) under future climate, meeting load targets for
all six climate models \parencite{Malone2020}, and watershed models estimate
$15$--$34\%$ reductions in tile nitrate yield
\parencite{David2015,Makhtoumi2025,Gupta2023}.

\subsection{Crop Rotation, Diversification, and Perennial Systems}

Cropping system exerts strong control on drainage nitrate, generally
exceeding tillage and often rivaling N rate. Continuous corn consistently lost
more nitrate than corn--soybean or perennial systems, leading authors to judge
it environmentally unsustainable \parencite{Weed1996,Kanwar1997}; rotations
produced greater yield and lower yield-scaled loss and a larger buffer above
the agronomic optimum \parencite{Zhao2016,Pasley2021}. Combining a shift to
corn--soybean with reduced rate and a cover crop lowered concentration from
$28$ to $8$~mg~L$^{-1}$ over 14 years \parencite{Ritter2015}, and more diverse
rotations delivered large reductions with maintained yield---a three-year
rotation averaged $4.2$ versus $8.0$~mg~L$^{-1}$ for a control, exceeding
$10$~mg~L$^{-1}$ in only $1\%$ of samples \parencite{Gentry2025,Ahmed2007b}.

Perennial and prairie systems are the most effective in-field option. Prairie
concentrations were $\le1.4$~mg~L$^{-1}$, $2$--$195$ times below row crops
\parencite{Daigh2015b}; restored prairie grass reduced losses to near rainfall
levels \parencite{HernandezRamirez2011}; and alfalfa and CRP plots had
$0.3$--$4$ versus $14$--$40$~mg~L$^{-1}$ for row crops, with four-year losses
of $\sim200$~kg~ha$^{-1}$ for corn versus $4$--$7$ for perennials
\parencite{Randall2001,Waring2024}. The benefit can be short-lived on
conversion back to row crops as RSN rises \parencite{Huggins2001}. Converting
corn--soybean area to alfalfa or wheatgrass reduced watershed losses
$78$--$81\%$ \parencite{Wilson2023}, and bioenergy perennials also reduce
loss \parencite{Ferin2023,Woodbury2018,Guo2018}. Diversified rotations, cover
crops, and perennials are recognized as complementary reducers
\parencite{Dinnes2002}.

\subsection{Tillage and Residue Management}

Tillage exerts weak and inconsistent control over drainage nitrate relative to
precipitation, cropping system, and N rate. Over 11 Minnesota years drain flow
was higher under no-till but nitrate losses only slightly higher under
conventional tillage, and moldboard-plow and no-till losses were similar
($43$ and $41$~kg~ha$^{-1}$), cropping system mattering far more
\parencite{RANDALL1995,Randall2001}. In Iowa no-till and ridge-till usually
had the lowest losses but rotation had greater influence
\parencite{Weed1996,Kanwar1997}, regional syntheses found no significant
tillage difference in N load \parencite{Zhao2016}, and conservation tillage
had the highest concentration ($16.1$~mg~L$^{-1}$) \parencite{Hertzberger2019,Wang2016b}.
Where reduced tillage builds soil health its nitrate effect remains equivocal:
no-till lowered loads only in some water years and reduced yields
\parencite{Masri2024}, though it raised soil organic carbon and lowered EONR
\parencite{Cassman2025,Waring2024}. Residue and stover management had limited
effect on nitrate: $50\%$ stover harvest reduced simulated loss only
$1.1$~kg~ha$^{-1}$ \parencite{Malone2019}, residue removal had little impact on
concentration or yield \parencite{Dougherty2020,Guo2018}, and continuous corn
with residue removal plus rye achieved losses similar to other row crops
\parencite{Daigh2015b,Farahani2020}.

\subsection{Soil Testing, Use Efficiency, Regional Aggregation, and Loss-Pathway Tradeoffs}

Watershed-scale soil-test-based programs confirm in-field mechanisms at
landscape scale. Adopting the late-spring nitrate test (LSNT) across a 400-ha
Iowa subbasin reduced concentration by at least $30\%$ via reduced fertilizer
use while matching yields \parencite{Jaynes2004,Karlen2005,Malone2010}, and
soil-test programs generally reduced RSN and loss versus yield-goal management
\parencite{Dinnes2002,Ritter2015}. Alternative farming practices reduced
drainage $41\%$ and concentration from $17.2$ to $8.2$~mg~L$^{-1}$, cutting
losses $59$--$62\%$ \parencite{Oquist2007}, and because $85\%$ of a Minnesota
watershed's load came from tile drainage, substantial perennialization was
needed for a $30$--$50\%$ reduction \parencite{Wilson2023,Makhtoumi2025}. Use
efficiency framing is increasingly emphasized: partial factor productivity was
highest for side-dress and lowest for very high fall N \parencite{Jaynes2015},
and EONR is proposed as a sustainability standard \parencite{Bandura2025};
adoption of cover crops, controlled drainage, and bioreactors remains
constrained by financial and social barriers \parencite{David2015}.

Scaling site results to the drained landscape indicates regionally significant
potential. For the five-state Corn Belt, winter rye was estimated to reduce
nitrate loss to the Mississippi by $\sim20\%$, with adoption on
$34$--$81\%$ of agricultural land and removal costs of
US\$$2.24$--$3.25$~kg$^{-1}$~N \parencite{Kladivko2014}; a statewide Illinois
analysis found rye reduced tile nitrate loss $\sim30\%$
\parencite{Gupta2023b}, and a five-state energy-cover-crop assessment
estimated a $216$~million~kg~yr$^{-1}$ reduction ($27\%$ of the assumed Gulf
load), with nine hotspot clusters on $19\%$ of the area contributing $48\%$ of
the reduction \parencite{Malone2023}.

Finally, mitigation practices redistribute N among loss pathways---leaching,
N$_2$O, and NH$_3$---so single-pathway evaluation can mislead. In-season N plus
rye cut leaching $37\%$ but raised annual N$_2$O $26\%$, leaving total N loss
and environmental-damage cost statistically unchanged \parencite{PrezaFontes2023}.
Yet split N plus rye had the lowest dissolved N$_2$O and nitrate loads, and no
consistent nitrate--N$_2$O tradeoff emerged across systems, showing stacked
practices need not trade one pathway for another
\parencite{Oliveira2025,OBrien2022}; ESN lowered N$_2$O, NH$_3$, and NO$_3^-$
simultaneously \parencite{Menegaz2025}. N$_2$O rose nonlinearly with rate above
a threshold near $120$~kg~N~ha$^{-1}$, and each unit of fertilizer-N produced
$0.01$~kg~N$_2$O-N alongside $0.12$~kg~NO$_3^-$-N loss, subsurface-drained
fields having the highest mean N$_2$O emission
\parencite{Wang2016b,Chatterjee2020b,Fontes2019}. The consistent conclusion
across scales is that combining two or more practices---rate optimization,
in-season timing, cover crops or diversified rotation, and drainage
management---reduces nitrate loss most cost-effectively, since in-field
management alone is often insufficient \parencite{Dinnes2002,Singh2023}.